

% Theorem environments
\usepackage{amsthm}  % For theorem-like environments
\theoremstyle{plain}
\newtheorem{theorem}{Theorem}
\newtheorem{lemma}[theorem]{Lemma}
\newtheorem{corollary}[theorem]{Corollary}
\newtheorem{proposition}[theorem]{Proposition}
\newtheorem{algorithm}[theorem]{Algorithm}
\newtheorem{assumption}[theorem]{Assumption}

\theoremstyle{definition}
\newtheorem{definition}{Definition}
\newtheorem{example}{Example}

\theoremstyle{remark}
\newtheorem{remark}{Remark}
    
\NewDocumentCommand{\br}{m}{\left( #1 \right)}
\NewDocumentCommand{\bs}{m}{\left[ #1 \right]}
\NewDocumentCommand{\bc}{m}{\left\{ #1 \right\}}

% Sets and number systems using dsfont (mathds)
\NewDocumentCommand{\R}{}{\mathds{R}}  % Real numbers
\RenewDocumentCommand{\H}{}{\mathds{H}}  % Hilbert space
\NewDocumentCommand{\N}{}{\mathds{N}}  % Natural numbers
\NewDocumentCommand{\Z}{}{\mathds{Z}}  % Integers
\NewDocumentCommand{\Q}{}{\mathds{Q}}  % Rational numbers
\NewDocumentCommand{\C}{}{\mathds{C}}  % Complex numbers
\NewDocumentCommand{\ind}{}{\mathds{1}}  % Complex numbers

% Mathematical operators and functions
\NewDocumentCommand{\defeq}{}{\coloneq} % Definition equal (
\NewDocumentCommand{\der}{}{\mathrm{d}}  % Differential oper
\NewDocumentCommand{\cder}{}{\partial}  % Differential oper
\NewDocumentCommand{\mul}{mm}{\langle #1, #2 \rangle}  % Inn
\NewDocumentCommand{\norm}{m}{\left\lVert #1 \right\rVert}  
\NewDocumentCommand{\abs}{m}{\left\lvert #1 \right\rvert}  %


\NewDocumentCommand{\maximize}{}{\text{maximize }}
\NewDocumentCommand{\minimize}{}{\text{minimize }}
\NewDocumentCommand{\subjectto}{}{\text{subject to }}
\NewDocumentCommand{\overall}{}{\text{over all }}


% Additional common mathematical abbreviations
\DeclareMathOperator*{\argmin}{arg\,min}
\DeclareMathOperator*{\argmax}{arg\,max}
\DeclareMathOperator{\dom}{dom}
\DeclareMathOperator{\epi}{epi}
\DeclareMathOperator{\prox}{prox}
\DeclareMathOperator{\conv}{conv}
\DeclareMathOperator{\tr}{tr} % Already using \operatorname,
\DeclareMathOperator{\rank}{rank} % Already using \operatorn

% More definitions as needed
\NewDocumentCommand{\grad}{}{\nabla}

\let\olddisplayopen\[
\let\olddisplayclose\]

% Redefine \[ and \]
\RenewDocumentCommand{\[}{}{%
  \begin{equation}\begin{aligned}%
}
\RenewDocumentCommand{\]}{}{%
  \end{aligned}\end{equation}%
}

